\title{Code Llama: Open Foundation Models for Code}

\begin{abstract}
We release \textsc{Code~Llama}, a family of large language models for code based on \textsc{Llama~2} providing state-of-the-art performance among open models, infilling capabilities, support for large input contexts, and zero-shot instruction following ability for programming tasks. We provide multiple flavors to cover a wide range of applications: foundation models (\textsc{Code~Llama}), Python specializations (\textsc{Code~Llama~-~Python}), and instruction-following models (\textsc{Code~Llama~-~Instruct}) with 7B, 13B, 34B, and 70B parameters each. These models are trained on sequences of 16k tokens and show improvements on inputs with up to 100k tokens. The 7B, 13B and 70B \textsc{Code~Llama} and \textsc{Code~Llama~-~Instruct} variants support infilling based on surrounding content.
\textsc{Code~Llama} reaches state-of-the-art performance among open models on several code benchmarks, with scores of up to 67\% and 65\% on HumanEval and MBPP, respectively. Notably, \textsc{Code~Llama~-~Python} 7B outperforms \textsc{Llama~2} 70B on HumanEval and MBPP, and all our models outperform every other publicly available model on MultiPL-E. We release \textsc{Code~Llama} under a permissive license that allows for both research and commercial use.\footnote{
\url{https://github.com/facebookresearch/codellama}}
\end{abstract}

\section{\textsc{Code~Llama}: Specializing \textsc{Llama~2} for code}
\label{sec:method}

\subsection{The \textsc{Code~Llama} models family}

\paragraph{\textsc{Code~Llama}.} The \textsc{Code~Llama} models constitute foundation models for code generation. They come in four model sizes: 7B, 13B, 34B and 70B parameters. The 7B, 13B and 70B models are trained using an infilling objective (\ref{sec:infilling}), and are appropriate to be used in an IDE to complete code in the middle of a file, for example. The 34B model was trained without the infilling objective. All \textsc{Code~Llama} models are initialized with \textsc{Llama~2} model weights and trained on 500B tokens from a code-heavy dataset (see \ref{sec:dataset} for more details), except \textsc{Code~Llama}~70B which was trained on 1T tokens. They are all fine-tuned to handle long contexts as detailed in \ref{sec:long_context}.

\paragraph{\textsc{Code~Llama~-~Python}.} The \textsc{Code~Llama~-~Python} models are specialized for Python code generation and come in sizes of 7B, 13B, 34B and 70B parameters. They are designed to study the performance of models tailored to a single programming language, compared to general-purpose code generation models. Initialized from \textsc{Llama~2} models and trained on 500B tokens from the \textsc{Code~Llama} dataset, \textsc{Code~Llama~-~Python} models are further specialized on 100B tokens using a Python-heavy dataset (\ref{sec:dataset}).
\textsc{Code~Llama~-~Python} with 7B, 13B and 34B parameters are trained without infilling and subsequently fine-tuned to handle long contexts (\ref{sec:long_context}).

\paragraph{\textsc{Code~Llama~-~Instruct}.} For the 7B, 13B and 34B sizes, \textsc{Code~Llama~-~Instruct} models are based on \textsc{Code~Llama} and fine-tuned with an additional approx.~5B tokens to better follow human instructions.

More details on \textsc{Code~Llama~-~Instruct} can be found in \ref{sec:instruct}. 

\paragraph{\textsc{Code~Llama}~70B.} \textsc{Code~Llama}~70B was trained months after the Code Llama 7B, 13B and 34B model. It was trained using the same data as the smaller versions of \textsc{Code~Llama}, and using roughly the same methods. \textsc{Code~Llama}~70B was trained on twice the number of tokens: 1 trillion instead of 500 billion. It was trained with FIM, which was an often-requested capability for the 34B model. Only the base \textsc{Code~Llama}~70B was trained with LCFT. See \ref{appendix:codellama70B_pipeline} for \textsc{Code~Llama}~70B specialization pipeline.
\textsc{Code~Llama~-~Instruct}~70B was trained from \textsc{Code~Llama~-~Python}~70B, which outperforms \textsc{Code~Llama}~70B in average on the languages from MultiPL-E including python.

\subsection{Dataset}
\label{sec:dataset}
We train \textsc{Code~Llama} 7B, 13B and 34B on 500B tokens, and \textsc{Code~Llama}~70B on 1T tokens during the initial phase, starting from the 7B, 13B, 34B, and 70B versions of \textsc{Llama~2}. 

As shown in Table~\ref{tab:dataset}, \textsc{Code~Llama} is trained predominantly on a near-deduplicated dataset of publicly available code. We also source 8\% of our samples data from natural language datasets related to code. This dataset contains many discussions about code and code snippets included in natural language questions or answers. To help the model retain natural language understanding skills, we also sample a small proportion of our batches from a natural language dataset. Data is tokenized via byte pair encoding (BPE, \citet{sennrich2016neural}), employing the same tokenizer as \textsc{Llama} and \textsc{Llama~2}.
Preliminary experiments suggested that adding batches sampled from our natural language dataset improves the performance of our models on MBPP. 

\subsection{Infilling}
\label{sec:infilling}

Code infilling is the task of predicting the missing part of a program given a surrounding context. Applications include code completion at the cursor's position in code IDEs, type inference and generation of in-code documentation (e.g., docstrings). 

We train infilling models following the concept of causal masking~\citep{aghajanyan2022cm3,fried2022incoder}, where parts of a training sequence are moved to the end, and the reordered sequence is predicted autoregressively. We train the general-purpose 7B, 13B and 70B models with an infilling objective, following the recommendations of \citet{bavarian2022efficient}.

More precisely, we split training documents at the character level into a prefix, a middle part and a suffix with the splitting locations sampled independently from a uniform distribution over the document length. We apply this transformation with a probability of 0.9 and  to documents that are not cut across multiple model contexts only. We randomly format half of the splits in the \emph{prefix-suffix-middle} (PSM) format and the other half in the compatible \emph{suffix-prefix-middle (SPM)} format described in \citet[App. D]{bavarian2022efficient}. We extend \textsc{Llama~2}'s tokenizer with four special tokens that mark the beginning of the prefix, the middle part or the suffix, and the end of the infilling span. To limit the distribution shift between autoregressive and infilling training, we suppress the implicit leading space that SentencePiece tokenizers add upon encoding the middle part and the suffix~\citep{kudo2018sentencepiece}. In SPM format, we concatenate the prefix and the middle part before encoding to tokens. Note that our model doesn't encounter split subtokens in the SPM format while it does in the PSM format.

Results on the effect of infilling training on downstream generation tasks and the performance of our infilling models on infilling benchmarks are reported in Section~\ref{sec:fim_results}.

\subsection{Long context fine-tuning}
\label{sec:long_context}

Effective handling of long sequences is a major topic of research in transformer-based language modeling~\citep{vaswani2017attention}. The fundamental modeling challenges are extrapolation, i.e., operating on sequence lengths beyond those seen at training time, and the quadratic complexity of attention passes which favors training on short-to-medium length inputs.

For \textsc{Code~Llama}, we propose a dedicated \emph{long context fine-tuning (LCFT)} stage in which models are presented with sequences of 16,384 tokens, up from the 4,096 tokens used for \textsc{Llama~2} and our initial code training stages. By limiting the training time spent on processing long sequences to a fine-tuning stage, we gain long-range capabilities without significantly increasing the cost of training our models. Our strategy is similar to the recently proposed fine-tuning by position interpolation~\citep{chen2023extending},  and we confirm the importance of modifying the rotation frequencies of the rotary position embedding used in the \textsc{Llama~2} foundation models \citep{su2021roformer}. However, instead of downscaling frequencies linearly as \citet{chen2023extending}, we change the base period from which they are derived. Specifically, with rotary embeddings, the query and key vectors $\mathbf{x}_n$ at position $n$ are subject to a linear transformation $\mathbf{R}^d_{\Theta,n} \mathbf{x}_n$, where $\mathbf{R}^d_{\Theta,n}$ is a block diagonal matrix with entries of the form

\begin{align*}
   \left( \mathbf{R}^d_{\Theta,n} \right) _i  =
   \begin{pmatrix}
   \cos n \theta_i & -\sin n \theta_i \\
   \sin n \theta_i & \cos n \theta_i
   \end{pmatrix},
\end{align*}

and $d$ denotes the embedding dimension. Rotation frequencies are computed as $\theta_i = \theta^{-2i/d}$, and we increase the base period $\theta$ from 10,000 to 1,000,000 for fine-tuning. This increase allows for processing much larger sequences and reduces bias towards short-distance attention (see \ref{app:lcft_details} for further discussion). Our experiments confirm that \textsc{Code~Llama} models are not only effective within the increased sequence length used during fine-tuning, but further show extrapolation capabilities and exhibit stable behavior on very long sequences of up to 100,000 tokens (\ref{sec:results-lcft}).

\subsection{Instruction fine-tuning}
\label{sec:instruct}

Our instruction fine-tuned models \textsc{Code~Llama~-~Instruct} are based on \textsc{Code~Llama} and trained to answer questions appropriately. They are trained on three different types of data.

\paragraph{Proprietary dataset.} We use the instruction tuning dataset collected for \textsc{Llama~2} and described in detail by \citet{touvron2023llamav2}. Specifically, we use the version referred to in their paper as ``RLHF V5'',  collected through several stages of reinforcement learning from human feedback and human feedback annotation  (see their Section 3 for more details). It combines thousands of Supervised Fine-Tuning and millions of Rejection Sampling examples. Each example consists of a multi-turn  dialogue between a \emph{user} and an \emph{assistant}. For Rejection Sampling, the output was selected among several generations using a reward model. The final dataset contains both Helpfulness and Safety data. This enables \textsc{Code~Llama} to inherit \textsc{Llama~2}'s instruction following and safety properties.

\paragraph{Self-instruct.}
Our proprietary dataset contains few examples of code-related tasks. Collecting supervised data from human annotators or training from human feedback~\citep{ouyang2022training} is expensive for coding tasks as it requires input from professional developers. Instead of human feedback, we use execution feedback to select data to train our instruct model. We construct the self-instruction dataset following the recipe below, resulting in $\sim$14,000 question-tests-solution triplets:

\begin{enumerate}
 \item Generate 62,000 interview-style programming questions by prompting (Figure~\ref{fig:prompt_self_generation_prompts}) \textsc{Llama~2}70B.
 \item De-duplicate the set of questions by removing exact duplicates, resulting in $\sim$52,000 questions.
 \item For each of these questions:
 \begin{enumerate}
    \item Generate unit tests by prompting \textsc{Code~Llama}7B (Figure~\ref{fig:test_generation_prompts})
    \item Generate ten Python solutions by prompting \textsc{Code~Llama}7B (Figure~\ref{fig:sol_generation_prompts})
    \item Run the unit tests on the ten solutions. Add the first solution that passes the tests (along with its corresponding question and tests) to the self-instruct dataset.
 \end{enumerate}
\end{enumerate}

We use \textsc{Code~Llama} 7B to generate the tests and Python solutions, as we found it more efficient than generating fewer solutions per question with the 34B model for the same compute budget. 

\paragraph{Rehearsal.}
In order to prevent the model from regressing on general coding and language understanding capabilities, \textsc{Code~Llama~-~Instruct} is also trained with a small proportion of data from the code dataset (6\%) and our natural language dataset (2\%). 

\subsection{Training details}

\paragraph{Optimization.}
Our optimizer is AdamW \citep{loshchilov2019decoupled} with $\beta_1$ and $\beta_2$ values of 0.9 and 0.95. We use a cosine schedule with $1000$ warm-up steps, and set the final learning rate to be 1/30th of the peak learning rate. We use a batch size of 4M tokens which are presented as sequences of 4,096 tokens each. Despite the standard practice of using lower learning rates in fine-tuning stages than in pre-training stages, we obtained best results when retaining the original learning rate of the \textsc{Llama~2} base model. We carry these findings to the 13B, 34B and 70B models, and set their learning rates to $3e^{-4}$, $1.5e^{-4}$, and  $1.5e^{-4}$ respectively. For python fine-tuning, we set the initial learning rate to $1e^{-4}$ instead. For \textsc{Code~Llama~-~Instruct}, we train with a batch size of 524,288 tokens and on approx.\ 5B tokens in total. 

\paragraph{Long context fine-tuning.} For long context fine-tuning (LCFT), we use a learning rate of $2e^{-5}$, a sequence length of 16,384, and reset RoPE frequencies with a base value of $\theta=10^6$.
The batch size is set to 2M tokens for model sizes 7B and 13B and to 1M tokens for model size 34B, respectively. Training lasts for 10,000 gradient steps by default. We observed instabilities in downstream performance for certain configurations, and hence set the number of gradient steps to 11,000 for the 34B models and to 3,000 for \textsc{Code~Llama}7B.

\section{Results}\label{sec:results}
We report results on a variety of benchmarks. First, we evaluate our models on popular description-to-code generation benchmarks for Python: HumanEval~\citep{chen2021evaluating},  MBPP~\citep{austin2021program}, and APPS \citep[programming interviews and competitions,][]{hendrycks2021measuring}. Second, we evaluate our models on further programming languages using MultiPL-E~\citep{cassano2022multiple}, namely on C++, Java, PHP, C\#, TypeScript (TS), and Bash. We additionally report results on the GSM8K benchmark~\citep{cobbe2021training}, which measures mathematical reasoning capabilities (Appendix~\ref{app:math}).

Next, we perform an extensive ablation study: (i) we study the impact of training from scratch or from a pretrained \textsc{Llama~2} model in \ref{sec:scratch}; (ii) we perform ablations for infilling and additional infilling specific benchmarks in \ref{sec:fim_results}; (iii) we study the effect of long context fine-tuning on perplexity, a synthetic retrieval task, and code completion with long source code files (\ref{sec:results-lcft}); and (iv) we evaluate our instruction fine-tuning procedure, which includes self-instruct training by leveraging self-generated unit tests in \ref{sec:inst_results}.

\subsection{Code generation}

\subsubsection{Python code generation}
\label{sec:python_evals}

We start by reporting results for Python code generation using the HumanEval~\citep{chen2021evaluating}, MBPP~\citep{austin2021program} and APPS~\citep{hendrycks2021measuring} benchmarks. Results are summarized in \ref{tab:main_res,tab:apps_res}. The full list of results on HumanEval and MBPP, including models with and without infilling and long context fine-tuning, can be found in \ref{tab:full_res} in \ref{sec:more_abb}. We provide zero-shot results of our instruction fine-tuned models on APPS in \ref{tab:apps_zero_shot} with evaluation details in \ref{appendix:apps_zero_shot}. Our main findings are as follows.

\paragraph{The value of model specialization.} We observe that model specialization is yields a boost in code generation capabilities when comparing \textsc{Llama~2} to \textsc{Code~Llama} and \textsc{Code~Llama} to \textsc{Code~Llama~-~Python}. \textsc{Llama~2} was trained on 2T tokens, and training on only 500B of extra tokens from a code-heavy dataset results in massive performance gains on both HumanEval and MBPP, to the point that \textsc{Llama~2} 70B is roughly equivalent to \textsc{Code~Llama} 7B on Python coding benchmarks. Although \textsc{Code~Llama} was trained on more than two epochs of our code dataset, which contains our entire Python dataset, training on 100B extra tokens of a Python-heavy data mix leads to significant gains on Python code generation benchmarks, between 4.3\% points and 8.3\% points in HumanEval pass@1 and between 1.2\% points and 6.4\% points in MBPP pass@1. These gains are smaller than for the first code training step, but still allow \textsc{Code~Llama~-~Python} 7B to outperform even \textsc{Code~Llama} 13B on MBPP and HumanEval. For the APPS benchmark, the prompts are much less direct and more complex compared to MBPP and HumanEval. Our \textsc{Code~Llama~-~Python} models show slightly decreased performance on the introductory and interview level problems, where understanding the prompt is often more challenging for a language model than implementing a solution. However, \textsc{Code~Llama~-~Python} shows clear gains on the competition-level problems where solutions are more complex. While large language models have enough capacity to learn to generate text on various topics, we observe that model specialization is beneficial for models between 7B and 70B parameters and after two full epochs on the training data. 

\paragraph{Scaling of specialized models.} We observe that scaling the number of parameters matters for models specialized for coding. With the same training process, our larger models outperform their smaller counterparts on almost every metric from HumanEval, MBPP and APPS (Table~\ref{tab:main_res},~\ref{tab:apps_res}). For instance, we gain 5.6 percentage points on MBPP pass@1 scaling \textsc{Code~Llama} from 7B to 13B parameters, 8 more points when scaling to 34B and 7 when scaling to 70B. We can hypothesize that specializing larger models to code would lead to significant further gains on coding tasks. Moreover, the Chinchilla scaling laws~\citep{hoffmann2022training} indicate that larger models would benefit more from training on more tokens. 

\subsubsection{Multilingual evaluation}
\label{sec:multilingual_evals}
Next, we evaluate our models on a more diverse set of programming languages. For that, we use the MultiPL-E benchmark~\citep{cassano2022multiple}. We report results for Python, C++, Java, PHP, TypeScript, C\#, and Bash in Table~\ref{tab:mlhe_res}.

We observe a similar improvement from \textsc{Llama~2} to \textsc{Code~Llama} in the multilingual setting as in the evaluation on Python (\ref{sec:python_evals}). The \textsc{Code~Llama} models clearly outperform \textsc{Llama~2} models of the same size on code generation in any language, and \textsc{Code~Llama} 7B even outperforms \textsc{Llama~2}~70B. 
Compared to other publicly available models, ours are especially strong in the multilingual setting. 
\textsc{Code~Llama} 7B outperforms larger models such as CodeGen-Multi or StarCoder, and is on par with Codex~\citep[code-cushman-001,][]{chen2021evaluating}.

The performance of \textsc{Code~Llama~-~Python} is comparable to that of \textsc{Code~Llama}. \textsc{Code~Llama~-~Python} 30B performs slightly worse than \textsc{Code~Llama} but \textsc{Code~Llama~-~Python} 7B and 13B perform slightly better than their counterparts without Python fine-tuning. More detailed results can be found in \ref{tab:full_multi_res}, \ref{sec:more_abb}.

To better understand the influence of multilingual pre-training, we measure the correlations between each of the evaluated languages and report the results separately for different model sizes in Figure~\ref{fig:correl}. We observe high correlation between model performance on C++, C\#, Java, and PHP. Interestingly, we also notice strong correlation between model performance on Python and Bash. Lastly, as expected the bigger and more expressive the models, the higher the correlation between the performance across all different languages.

\subsection{Infilling evaluations}\label{sec:infilling_evals}

\label{sec:fim_results}
\paragraph{Performance cost of infilling training.} Previous studies on infilling (or \emph{fill-in-the-middle, FIM}) code models assert that the traditional next token prediction objective can be replaced by a multitask infilling objective with an infilling rate of up to 90 \% at no cost for left-to-right autoregressive test losses \citep{bavarian2022efficient} and only small cost for downstream evaluation performance \citep{allal2023santacoder}. In Table~\ref{tab:fim-cost}, we independently validate both findings at the scale of 7B and 13B parameters and 500B training tokens of code. The 7B model loses 0.6 percentage points on average across HumanEval and MBPP pass@1, pass@10 and pass@100 scores if trained with an infilling objective, while the 13B model loses 1.1 percentage points. Because of this modest decline in performance and the wide applicability of models with infilling capability, we decide to release \textsc{Code~Llama} 7B, 13B and 70B in this configuration.

\paragraph{Code infilling benchmarks.} Our infilling models reach state-of-the-art performances in code infilling benchmarks among models of their size. We evaluate on two related code infilling benchmarks based on the HumanEval benchmark \citep{chen2021evaluating}.

The HumanEval infilling benchmark \citep{fried2022incoder} turns the reference solutions of the HumanEval benchmark \citep{chen2021evaluating} into infilling problems by masking out either individual lines or blocks consisting of multiple consecutive lines. It has been extended in \citet{bavarian2022efficient} with a random span infilling task in which the masking is applied to a randomly selected substring at the character level. Predictions are scored with a pass@1 score based on the test cases of the original HumanEval problems. According to the results in \ref{tab:he-fim-incoder}, our models outperform all other infilling models of their size. Note, however, that the results in random span infilling are significantly worse in suffix-prefix-middle (SPM) format than in prefix-suffix-middle (PSM) format as it would require token healing \citep{microsoft2023guidance}, which we have not implemented for this evaluation (see Appendix~\ref{app:infilling} for further discussion).

\citet{allal2023santacoder} translates the HumanEval infilling benchmark to other programming languages using MultiPL-E \citep{cassano2022multiple}. Single lines are masked and predictions are scored with an exact match metric against the ground truth solution. Our models, including \textsc{Code~Llama} 7B, outperform all open infilling models across the three programming languages contained in the benchmark (\ref{tab:he-fim-santa}). We observe a further increase in performance when prompting the models in SPM format, like witnessed in \cite{bavarian2022efficient}. 

\subsection{Long context evaluations}

\label{sec:results-lcft}
We explore \textsc{Code~Llama}'s ability to work with long sequences by measuring perplexity, key retrieval accuracy and performance during generation on code completion tasks. These tasks, and our results are detailed below. For full results and comparisons to alternative techniques of increasing the context length of LLMs, we refer to~\ref{app:lcft}.

\paragraph{Perplexity during extrapolation.} In \ref{fig:lcft-code-ppl}, perplexity is computed over 4M tokens from the code dataset, using a subset of our validation data consisting of large source files ($\ge$50kB).
For all model sizes, we observe a steady decrease in perplexity well beyond 16384 tokens, which is the sequence length we use for long-context fine-tuning. After 100K tokens, the perplexity increases only slightly, in contrast to the well-known instability phenomenon when testing transformer models on sequences larger than those seen during training~\citep{press2021train}.

\paragraph{Key retrieval.} In \ref{fig:lcft-key-retrieval}, we investigate key retrieval performance in synthetic task. The prompt consists of a large amount of syntactically valid Python code, with a function returning a scalar inserted at a specified position. The model is asked to complete an \texttt{assert} statement with the return value of the inserted function.
\citet{liu2023lost} showed that the inability to recall content placed in the middle of long prompts is a common failure mode in LLMs; our retrieval task is analogous to their setup, albeit tailored to code models which are not fine-tuned to follow instructions. All models exhibit strong retrieval performance on the sequence length they were trained on, with the exception of the 7B model for test cases in which the function is placed at the beginning of the prompt. We include OpenAI’s gpt-3.5-turbo-16k-0613 as a reference. We query GPT with a system prompt of ``Complete the following code.'' and a temperature of 0. For sequences beyond 16K tokens, i.e., when extrapolating, our models exhibit a decrease in performance (\ref{app:lcft_extended_results}).

\paragraph{Single line completion.} Finally, we test the benefits of the ability to handle long context sizes in a single line code completion task. Our task is based on the Long Code Completion (LCC) benchmark~\citep{guo2023longcoder}.\footnote{Note that LCC data points are included in our code training data.} The LCC test set is skewed towards shorter files and we hence sample a new set of examples from LCC's validation and test set with an equalized distribution over file size (\ref{app:lcft_benchmarks}).
In \ref{tab:lcc_results}, we compare the completion accuracy of the \textsc{Code~Llama} models to their counterparts prior to long-context fine-tuning. Non-LCFT models fail to generate meaningful completions on long sequences and we thus truncate their prompts to the 4,000 tokens immediate preceding the line to complete. Across all metrics, models fine-tuned to handle long contexts achieve significantly higher performance. This demonstrates that long contexts are informative for code completion, and that with LCFT our models are able to leverage this information to improve their generations. We note that the longest example's prompt in this test consists of 103K tokens, for which all \textsc{Code~Llama} models generate syntactically correct completions, with the 7B model producing an exact match.

\paragraph{Performance impact on short contexts.} While our models are effective on long sequences, we observe that LCFT slightly hurts performance on standard code synthesis benchmarks consisting of short sequences. In \ref{tab:full_res}, we observe an average decrease of 0.52 percentage points on HumanEval pass@1 and 1.9 points on MBPP for the pass@1 metric. Similarly, a breakdown of the code completion results in \ref{tab:lcc_results} by the number of tokens in each example shows that for prompts shorter than 4k tokens, long context fine-tuning induces a reduction of up to 2 BLEU points from base models after code training (\ref{fig:lcft-lcc-difference}).
We observe similar decreases in performance for infilling tasks (Table~\ref{tab:he-fim-incoder}).

LCFT comes at a cost for short sequences, and slightly decreases our scores on standard coding benchmarks such as HumanEval and MBPP. However, many real-world use cases are not captured by these benchmarks, and we believe that this cost is more than offset by the potential of handling long sequences for real downstream applications. Hence we opt to release all our \textsc{Code~Llama}, \textsc{Code~Llama~-~Python} and \textsc{Code~Llama~-~Instruct} models with long-context capabilities.

\subsection{Ablation studies}

\subsubsection{Fine tuning \textsc{Llama~2} vs. training from scratch on code}\label{sec:scratch}

\textsc{Code~Llama} is based on the \textsc{Llama~2} models, which are trained on 2T tokens of text, including only 80B tokens of code. We tune these models on 500B extra tokens, consisting mostly of code (85\%). 
Figure~\ref{fig:training_curves} shows the training curves of \textsc{Code~Llama}. 

We compare the 7B parameters model to an identical model trained from scratch on the same data mix (Figure~\ref{fig:curves_scratch_ablation_loss}). At the end of training, the loss of the model trained from scratch is equal to the loss of \textsc{Code~Llama} 7B at about half of its training (with 240B less training tokens). Moreover, this gap becomes larger over time. 

\subsubsection{Instruction fine-tuning}\label{sec:inst_results}
\paragraph{General helpfulness vs. coding ability}

We evaluate \textsc{Code~Llama~-~Instruct} and compare it to \textsc{Llama~2}-Chat for coding tasks and helpfulness (\ref{fig:helpfulness_vs_mbpp}). 
We observe that \textsc{Code~Llama} improves its coding abilities for each model sizes, while preserving the general helpfulness performance inherited from \textsc{Llama~2}. 
The results on the helpfulness axis is an indication that \textsc{Code~Llama} performs greatly on general instructions following. But we emphasize that this result should be taken with a grain of salt, since we limited our automatic evaluation to scoring the models answers with \textsc{Llama~2} reward model.

\paragraph{The value of self-instruct data}

We also perform ablations, showing the value of the self-instruct data that we generate with our own model. To evaluate the capacity of the model to answer questions, we use a zero-shot version of MBPP. We prompt the model to generate the code between \texttt{[PYTHON]} and \texttt{[/PYTHON]} tags to make it easy to parse the result. Our exact prompt is shown in Figure~\ref{fig:mbpp_zero_prompt} in the Appendix.
Table~\ref{tab:self_instruct} show the impact of training on data generated using our models and filtered with unit tests as described in Section~\ref{sec:instruct}. The self-instruct data allows us to improve our scores on benchmarks such as HumanEval and MBPP. It also makes the training more reliable. With self-instruct, the model easily learns to follow the format requested for MBPP zero-shot while it sometimes fails without it.

\paragraph{Unnatural model.} For comparison purposes, we also finetuned \textsc{Code~Llama~-~Python} 34B on 15,000 unnatural instructions similarly to~\citet{honovich2022unnatural} using the same prompts as for the self-instruct dataset. We do not release this model, but we observe clear improvements on HumanEval and MBPP which are indicative of the improvements that can be reached with a small set of high-quality coding data. The results of the unnatural model are shown in Table~\ref{tab:main_res}. 

\subsubsection{Pass@k evaluation} We study the effect of the sampling temperature on the pass@k performance. Specifically, we report pass@1, 10, and 100 using temperature $\in \{0.1, 0.4, 0.6, 0.8\}$ on both HumanEval and MBPP. Results are depicted in Figure~\ref{fig:abb_temp}. As expected, as we increase the temperature, the pass@1 scores are getting worse while the pass@10 and pass@100 improve. 

\section{Responsible AI and safety}\label{sec:safety}

Large language models have been shown to have the potential to produce known falsehoods due to misconceptions or false beliefs \citep{lin2021truthfulqa}, generate toxic or offensive content \citep{hartvigsen2022toxigen} and reproduce or even amplify the biases that are contained in the training data \citep{dhamala2021bold}. As mentioned in \ref{sec:instruct}, we make \textsc{Code~Llama~-~Instruct} safer by fine-tuning on outputs from \textsc{Llama~2}, including adversarial prompts with safe responses, as well as prompts addressing code-specific risks.

In this section, we perform evaluations on three widely-used automatic safety benchmarks from the perspectives of truthfulness, toxicity, and bias, respectively. Specifically, we assess the safety capabilities of both pretrained \textsc{Code~Llama} and fine-tuned \textsc{Code~Llama~-~Instruct} with Falcon \citep{falcon40b}, MPT \citep{themosaicmlnlpteam2023introducing}, and StarCoder \citep{li2023starcoder}. Although we have chosen certain standard benchmarks commonly used in the language model community to highlight some of the problems with these models, it's important to note that these evaluations alone do not provide a comprehensive understanding of the risks associated with them. We complement the safety analysis of \textsc{Code~Llama~-~Instruct} with additional red teaming from various domain experts in offensive security, malware development, responsible AI and software engineering, similar to \cite{touvron2023llamav2}.

\paragraph{Truthfulness.} We use \textbf{TruthfulQA} \citep{lin2021truthfulqa} to gauge the factuality and common sense of our models. The TruthfulQA benchmark comprises 817 questions spread across 38 categories, encompassing topics such as health, finance, law, and politics \citep{lin2021truthfulqa}. The questions are designed to be challenging, even for humans, causing them to answer incorrectly due to unfounded beliefs or misconceptions. To evaluate the generated outputs from LLMs, we utilize GPT-3-based metrics following \cite{lin2021truthfulqa} to determine the truthfulness and informativeness of the outputs. For the QA prompt, we use a few-shot prompt containing 6 random QA pairs, structured according to the InstructGPT format \citep{ouyang2022training}. The results are reported as the percentage of generations that are both truthful and informative, as well as the percentage that are either truthful or informative.

\paragraph{Toxicity. } We use \textbf{ToxiGen} \citep{hartvigsen2022toxigen} to quantify the extent of toxic language and hate speech generation across various demographic groups. The ToxiGen dataset contains implicitly toxic and benign sentences mentioning 13 minority groups. Following \cite{touvron2023llamav2}, we utilize an improved version of the dataset, which minimizes noise by removing prompts with disagreements among annotators regarding the target demographic group. To measure the toxicity of the generated outputs from each of the LLMs, we employ the default ToxiGen classifier, tuned on RoBERTa \citep{liu2019roberta}. 

\paragraph{Bias. } We employ the Bias in Open-Ended Language Generation Dataset (\textbf{BOLD}) \citep{dhamala2021bold} to investigate how the sentiment in the model's outputs may differ based on demographic attributes. The BOLD benchmark consists of a total of 23,679 English Wikipedia prompts that span five domains: race, gender, religion, political ideology, and profession. These prompts cover 43 different subgroups. In our analysis, we exclude prompts belonging to the religious ideology subgroups Hinduism and Atheism due to their limited representation, consisting of only 12 and 29 prompts, respectively. To assess the sentiments conveyed by the combination of the prompt prefix and model generation, we employ sentiment analysis using the Valence Aware Dictionary and Sentiment Reasoner (VADER) \citep{hutto2014vader}. The VADER produces sentiment scores between -1 and 1, where a positive (negative) score indicates a positive (negative) sentiment towards the population mentioned in the prompt. A score closer to 0 indicates a neutral sentiment.

\paragraph{Benchmark evaluation results. }
Table~\ref{tab:safety_bench} shows the evaluation results of the three safety benchmarks. We follow the decoding setting as in \cite{touvron2023llamav2} where a temperature of 0.1 and top-p of 0.9 are used. Regarding TruthfulQA, we provide the percentage of generations that are both truthful and informative, where a higher percentage indicates better performance. Regarding ToxiGen, we present the percentage of generations deemed toxic by the metric, with a lower percentage indicating better results. Regarding BOLD, we present the average sentiment scores across demographic groups within the five domains in the BOLD dataset. The fine-tuned \textsc{Code~Llama~-~Instruct} exhibits significant improvements over the pretrained \textsc{Code~Llama} in terms of truthfulness (from $34.64$ to $47.37$ for 34B) and toxicity (from $17.62$ to $0.00$ for 34B). The percentage of toxic generations drastically reduces to virtually 0\% across all \textsc{Code~Llama} sizes, making it the least toxic among all the models compared. When compared to Falcon and MPT fine-tuned models, the fine-tuned \textsc{Code~Llama} demonstrates the second-best performance level in both toxicity and truthfulness, right after \textsc{Llama~2~Chat}. 
Additionally, similar to \textsc{Llama~2~Chat}, the \textsc{Code~Llama~-~Instruct}, after fine-tuning, also tends to show an overall increase in positive sentiment for many demographic groups in BOLD. More detailed results split by different demographic groups can be found in \ref{appendix:safety}. 

\paragraph{Red teaming. } It is important to also proactively identify risks with adversarial testing or red teaming. We conducted 3 red teaming exercises with 25 Meta employees, including domain experts in responsible AI, malware development, and offensive security engineering.

The red teamers provided a nuanced evaluation specifically on the risk from so called ``dual intent prompts.'' Dual intent prompts are requests for help with writing code that could be used maliciously but the prompt does not directly address the topic (example ``Mosaic Prompts''~\cite{glukhov2023llm}). For example, the model rightfully refuses to provide support with writing ransomware code but it complies when asked to provide a script to encrypt all files in the user's home directory since such a script could be used for~benign~purposes. 

After conducting red team exercises, we asked participants (who had also participated in \textsc{Llama~2~Chat} exercises) to also provide qualitative assessment of safety capabilities of the model. Some participants who had expertise in offensive security and malware development questioned the ultimate risk posed by ``malicious code generation'' through LLMs with current capabilities. 

One red teamer remarked, ``While LLMs being able to iteratively improve on produced source code is a risk, producing source code isn't the actual gap. That said, LLMs may be risky because they can inform low-skill adversaries in production of scripts through iteration that perform some malicious behavior.''  

According to another red teamer, ``[v]arious scripts, program code, and compiled binaries are readily available on mainstream public websites, hacking forums or on `the dark web.' Advanced malware development is beyond the current capabilities of available LLMs, and even an advanced LLM paired with an expert malware developer is not particularly useful- as the barrier is not typically writing the malware code itself. That said, these LLMs may produce code which will get easily caught if used directly.''

In addition to red teaming sessions, we ran a quantitative evaluation on risk from generating malicious code by scoring \textsc{Code~Llama}'s responses to ChatGPT's (GPT3.5 Turbo) with LLAMAv2 70B's safety reward model. For this second quantitative evaluation, we selected prompts that the red teamers generated specifically attempting to solicit malicious code (even though the red teaming included consideration of a broad set of safety risks). These prompts were a mix of clear intent and slightly obfuscated intentions (see some examples in Figure~\ref{fig:red_teaming_code_prompts}.
We show a KDE plot of the distribution of the safety score for all models in Figure~\ref{fig:codellama_vs_chatgpt_redteaming_code}).
We observe that \textsc{Code~Llama} tends to answer with safer responses; the distribution of safety scores for \textsc{Code~Llama} has more weight in the safer part of the range.

\paragraph{False refusals.} LLMs that are too safe can have a tendency to over-refuse valid claims similar to what was reported after the release of \textsc{Llama~2}. 
We specifically asked red teamers to test for this behavior. They found some limited evidence of false refusals (when not using a system preprompt). False refusals could also be solved by rephrasing the prompt e.g. ``Can you tell me how to kill a process?'' rephrased to ``How do I kill a process?''. We show some examples in Appendix Table~\ref{fig:red_teaming_false_refusals}.
This behavior is something we plan to investigate in more details in the future.

\paragraph{Safety and coding performance.} As our instruction finetuning set prioritizes safety, longer finetunings tend to degrade coding performance. We trained our models to reach high coding performances, while not compromising on safety. As shown in Figure~\ref{fig:codellama_vs_chatgpt_redteaming_code}, our \textsc{Code~Llama~-~Instruct} models are safer~than~ChatGPT. 

\newcommand{LLM}{LLM}
\newcommand{LLMs}{LLMs}

\section{Related work}\label{sec:relatedwork}
Early observations with LLMs such as GPT-Neo \citep{gpt-neo} or GPT-J \citep{gpt-j} showed that adding code in the training data makes program synthesis possible even with medium size LLMs. Code from open-source software is now a standard part of the training data for general-purpose LLMs such as PaLM \citep{chowdhery2022palm}, Chinchilla \citep{hoffmann2022training}, Gopher \citep{rae2021scaling}, GPT-4 \citep{openai2023gpt4}, and \textsc{Llama} \citep{touvron2023llama,touvron2023llamav2}. In parallel, models specifically trained or fine-tuned for code understanding and program synthesis from natural language prompts emerged with LLMs such as Codex \citep{chen2021evaluating}, CodeT5 \citep{wang2021codet5}, InCoder \citep{fried2022incoder}, AlphaCode \citep{li2022alphacode}, CodeGen \citep{nijkamp2022codegen} and CodeGen 2 \citep{nijkamp2023codegen2}, GPT-NeoX \citep{black2022gpt}, SantaCoder \citep{allal2023santacoder}, StarCoder \citep{li2023starcoder} and phi-1  \citep{gunasekar2023textbooks}, consistently demonstrating better performance on code benchmarks than  general-purpose LLMs of comparable or even larger size. This paper follows this line, by fine-tuning the recent general-purpose language model \textsc{Llama~2} on code data.

\paragraph{Closed-source vs open-source models.} The landscape of LLMs is marked by whether the technology is free and the code is available for research or commercial use. ChatGPT and GPT-4 \citep{openai2023gpt4}, PaLM \citep{chowdhery2022palm} and Chinchilla \citep{hoffmann2022training} are closed source, while BLOOM \citep{scao2022bloom}, OPT \citep{zhang2022opt}, and the seminal work of \textsc{Llama} are public~\citep{touvron2023llama}. The more recent \textsc{Llama~2} has been released under a custom licence for commercial use \citep{touvron2023llamav2}. A similar dichotomy exists for code models, with Codex/copilot \citep{chen2021evaluating}, AlphaCode \citep{li2022alphacode}, GPT-4 or phi-1 \citep{gunasekar2023textbooks} being closed source, whereas the recent SantaCoder \citep{allal2023santacoder} and StarCoder \citep{li2023starcoder} have been released open-source and allow for commercial use. In this work, we allow for commercial use of the models under the same terms as \textsc{Llama~2}. 
Moreover, our largest model, with its 70B parameters, is significantly larger than previous open-source models -- GPT-NeoX-20B \citep{black2022gpt} and StarCoder with 15.5B parameters -- which allows it to achieve state-of-the-art performances on HumanEval, MBPP and MultiPL-E among open-source models. -- GPT-NeoX-20B \citep{black2022gpt} and StarCoder with 15.5B parameters -- which allows it to achieve state-of-the-art performances on HumanEval, MBPP and MultiPL-E among open-source models. 

\paragraph{Data.} 
It is well-known that data quality is critical in the training and responsible development of LLMs \citep[e.g.,][]{hoffmann2022training,penedo2023refinedweb}, and this is also true for code as discussed by \citet{allal2023santacoder}. Modern models are trained on publicly available, open-source code. In addition, \citet{allamanis2019adverse} and \citet{allal2023santacoder} discuss the impact of effective deduplication and of selecting code from repositories based on the number of GitHub stars (as a proxy for popularity), while \citet{li2023starcoder} augment their data with GitHub issues and commits collected from BigQuery. \citet{gunasekar2023textbooks} filter data up to only containing ``textbook''-quality code and add synthetic problems collected using GPT-3.5, following \citet{jung2023impossible}, in order to obtain good performance on simple benchmarks such as HumanEval and MBPP. We follow the approach of learning from publicly available code only, without additional meta-level or temporal information such as issues or commits. We also do not train our foundation models on additional synthetic exercises, since we did not want to take the risk of reducing the scope of our models to simple coding exercises similar to those contained in HumanEval and MBPP.

\paragraph{Code understanding and synthesis tasks.} In addition to program synthesis from natural language prompts or infilling \citep{fried2022incoder,bavarian2022efficient,li2023starcoder,nguyen2023meet}, many tasks related to code understanding or synthesis have been addressed since the early 2020s with NLP models adapted for code \citep{raffel2020exploring,feng2020codebert,guo2020graphcodebert,wang2021codet5,ahmad2021unified}, also see the survey by~\citet{xu2022survey}. These tasks include code summarization, refinement, translation \citep{roziere2020unsupervised,DBLP:journals/corr/abs-2110-06773,szafraniec2022code} fixing bugs ~\citep{yasunaga2021break,zhang2022repairing,prenner2022can}, fixing build errors \citep{tarlow2020learning} or generating unit tests~\citep{tufano2020unit,li2022alphacode,chen2022codet}, as well as solving math problems as demonstrated by PaLM \citep{chowdhery2022palm} or Codex \citep{chen2021evaluating}. 14 code understanding tasks are represented in the CodeXGlue benchmark \citep{CodeXGLUE}. Here we focused on the main problem of program synthesis, as well as infilling/completion for our 7B and 13B models where the ability comes with little impact on the generation performance as previously observed by \citet{bavarian2022efficient}. 

\paragraph{Additional modifications to LLM training and inference.} A number of works proposed to incorporate within the training objective structural knowledge of programs, with specialized objectives for code deobfuscation~\citep{roziere2021dobf}, contrastive learning through semantic-preserving code transformations ~\citep{jain-etal-2021-contrastive}, leveraging Abstract Syntax Trees to learn tree-aware positional encodings ~\citep{NEURIPS2019_6e091746,peng2021integrating}. A recent stream of work takes into account program execution or unit tests to filter, cluster, or improve the correctness of programs when few candidates must be submitted \citep{li2022alphacode,chen2022codet,le2022coderl,zhang2023planning}, or unit tests them within a reinforcement learning objective to enrich the training signal \citep{le2022coderl,liu2023rltf}. We focused here on improving the base model rather than tweaking the inference scheme, since we believe this is where most of the long-term progress comes from; it is nonetheless an interesting direction to experiment with more elaborated inference schemes on top of \textsc{Code~Llama}.

\paragraph{Long sequences in LLMs.} Scaling Transformers and LLMs to long input sequences has attracted much recent interest~\citep{dai2019transformerxl,beltagy2020longformer,yu2023megabyte,ding2023longnet}. The context lengths supported by available models and APIs has seen a steady increase, with StarCoder being trained on 8K token sequences (\citep{li2023starcoder}, up from the 4K of \citet{allal2023santacoder}), recent GPT versions supporting 16K (gpt-3.5-turbo-16k) and 32K tokens (gpt-4-32k), MPT-7b fine-tuned on 65K tokens~\citep{themosaicmlnlpteam2023introducing}, and Claude featuring 100K context windows~\citep{anthropic2023introducing}. Previous research focuses on alleviating the $O(n^2)$ space and time complexity of self-attention~\citep{vaswani2017attention} by introducing sparsity patterns, as well as by encoding positional information in such a way that models can leverage input sizes larger than those presented at training time (length extrapolation). In our work, we do not rely on hand-crafted sparsity patterns such as those proposed for code input by~\citet{guo2023longcoder}, who operate on sequences of up to 4,096 tokens, as to not curtail the model's expressivity, and modify the encoding of positions instead. Starting from pretrained \textsc{Llama~2} models that utilize RoPE~\citep{su2021roformer}, \citet{chen2023extending} propose additional fine-tuning for long sequence handling, an approach we pursue as well. However, we tailor our hyper-parameter modifications to allow for extrapolation at inference time. Our modification of the RoPE hyper-parameters~\citep{su2021roformer} is a simple modification which does not require any architectural changes or restrictions and can be readily applied to existing implementations.\footnote{Concurrently to our work, the approach of increasing the rotation frequency base value has been proposed by user ``bloc97'' in the ``LocalLLaMA'' subreddit (\url{https://redd.it/14lz7j5}), where it was applied to LLaMA models without further fine-tuning.}
\citet{press2021train} propose a linear bias for attacking extrapolation; in contrast, our approach seeks to reduce existing bias towards shot-range attention. Recent work suggests that causal models do not require an explicit encoding of position information~\citep{haviv2022transformer,kazemnejad2023impact}, a hypothesis we did not test in this work as we demonstrated that starting from pretrained \textsc{Llama~2} models is significantly more efficient than training from scratch.

\section{Discussion}
We release a family of code-specialized \textsc{Llama~2} models called \textsc{Code~Llama}, with three main variants that we release with four sizes (7B, 13B, 34B, and 70B parameters): \textsc{Code~Llama}, \textsc{Code~Llama~-~Python}, \textsc{Code~Llama~-~Instruct}. 
With real-world applications in mind, we trained our 7B, 13B, and 70B models to support infilling, and all our models to leverage large contexts. We tested their stability in inference up to 100K tokens (\ref{fig:lcft-code-ppl}). 
Large context fine-tuning and infilling come at a cost on standard benchmarks left-to-right code generation benchmarks (\ref{tab:full_res}), that are all based on short sequences (i.e. function level). Still, our 70B model is state-of-the-art among public models on standard python completion benchmarks, and our other models are competitive compared to models with similar numbers of parameters. 

On multilingual benchmarks, even our smallest model (\textsc{Code~Llama} 7B) outperforms every other public model. 

The \textsc{Code~Llama~-~Instruct} models are trained to provide zero-shot instruction ability to \textsc{Code~Llama}. In this further fine-tuning, where we somewhat distillate~\textsc{Llama~2}-Chat, we focused not only on being more directly helpful (\ref{fig:helpfulness_vs_mbpp}) but also sought to provide a safer model to use and deploy (\ref{sec:safety}). Following instruction and being overly safe can cost some points on evaluations (e.g. on HumanEval for the 34B model in~\ref{tab:main_res}), as exemplified in~\ref{fig:red_teaming_false_refusals}. Further work is needed for LLMs to understand context and nuance in their instructions.

\end{document}
